\documentclass[a4paper, 14pt]{extarticle}

% Поля
%--------------------------------------
\usepackage{geometry}
\geometry{a4paper,tmargin=2cm,bmargin=2cm,lmargin=3cm,rmargin=1cm}
%--------------------------------------


%Russian-specific packages
%--------------------------------------
\usepackage[T2A]{fontenc}
\usepackage[utf8]{inputenc} 
\usepackage[english, main=russian]{babel}
%--------------------------------------

\usepackage{textcomp}

% Красная строка
%--------------------------------------
\usepackage{indentfirst}               
%--------------------------------------             


%Graphics
%--------------------------------------
\usepackage{graphicx}
\graphicspath{ {./images/} }
\usepackage{wrapfig}
%--------------------------------------

% Полуторный интервал
%--------------------------------------
\linespread{1.3}                    
%--------------------------------------

%Выравнивание и переносы
%--------------------------------------
% Избавляемся от переполнений
\sloppy
% Запрещаем разрыв страницы после первой строки абзаца
\clubpenalty=10000
% Запрещаем разрыв страницы после последней строки абзаца
\widowpenalty=10000
%--------------------------------------

%Списки
\usepackage{enumitem}

%Подписи
\usepackage{caption} 

%Гиперссылки
\usepackage{hyperref}

\hypersetup {
	unicode=true
}

%Рисунки
%--------------------------------------
\DeclareCaptionLabelSeparator*{emdash}{~--- }
\captionsetup[figure]{labelsep=emdash,font=onehalfspacing,position=bottom}
%--------------------------------------

\usepackage{tempora}
\usepackage{amsmath}
\usepackage{color}
\usepackage{listings}
\lstset{
  belowcaptionskip=1\baselineskip,
  breaklines=true,
  frame=L,
  xleftmargin=\parindent,
  language=Python,
  showstringspaces=false,
  basicstyle=\footnotesize\ttfamily,
  keywordstyle=\bfseries\color{blue},
  commentstyle=\itshape\color{purple},
  identifierstyle=\color{black},
  stringstyle=\color{red},
}

%--------------------------------------
%			НАЧАЛО ДОКУМЕНТА
%--------------------------------------

\begin{document}

%--------------------------------------
%			ТИТУЛЬНЫЙ ЛИСТ
%--------------------------------------
\begin{titlepage}
\thispagestyle{empty}
\newpage


%Шапка титульного листа
%--------------------------------------
\vspace*{-60pt}
\hspace{-65pt}
\begin{minipage}{0.3\textwidth}
\hspace*{-20pt}\centering

\includegraphics[width=\textwidth]{emblem}
\end{minipage}
\begin{minipage}{0.67\textwidth}\small \textbf{
\vspace*{-0.7ex}
\hspace*{-6pt}\centerline{Министерство науки и высшего образования Российской Федерации}
\vspace*{-0.7ex}
\centerline{Федеральное государственное бюджетное образовательное учреждение }
\vspace*{-0.7ex}
\centerline{высшего образования}
\vspace*{-0.7ex}
\centerline{<<Московский государственный технический университет}
\vspace*{-0.7ex}
\centerline{имени Н.Э. Баумана}
\vspace*{-0.7ex}
\centerline{(национальный исследовательский университет)>>}
\vspace*{-0.7ex}
\centerline{(МГТУ им. Н.Э. Баумана)}}
\end{minipage}
%--------------------------------------

%Полосы
%--------------------------------------
\vspace{-25pt}
\hspace{-35pt}\rule{\textwidth}{2.3pt}

\vspace*{-20.3pt}
\hspace{-35pt}\rule{\textwidth}{0.4pt}
%--------------------------------------

\vspace{1.5ex}
\hspace{-35pt} \noindent \small ФАКУЛЬТЕТ\hspace{80pt} <<Информатика и системы управления>>

\vspace*{-16pt}
\hspace{47pt}\rule{0.83\textwidth}{0.4pt}

\vspace{0.5ex}
\hspace{-35pt} \noindent \small КАФЕДРА\hspace{50pt} <<Теоретическая информатика и компьютерные технологии>>

\vspace*{-16pt}
\hspace{30pt}\rule{0.866\textwidth}{0.4pt}
  
\vspace{11em}

\begin{center}
\Large {\bf Лабораторная работа № 2} \\ 
\large {\bf<<Реализация метода Гаусса c перестановками>>} \\ 
\large {\bf по курсу <<Численные методы линейной алгебры>>} \\ 
\end{center}\normalsize

\vspace{8em}


\begin{flushright}
  {Студент группы ИУ9-71Б Лысенко В. В.\hspace*{15pt} \\
  \vspace{2ex}
  Преподаватель Посевин Д. П.\hspace*{15pt}}
\end{flushright}

\bigskip

\vfill
 

\begin{center}
\textsl{Москва 2024}
\end{center}
\end{titlepage}
%--------------------------------------
%		КОНЕЦ ТИТУЛЬНОГО ЛИСТА
%--------------------------------------

\renewcommand{\ttdefault}{pcr}

\setlength{\tabcolsep}{3pt}
\newpage
\setcounter{page}{2}

\section{Цель работы}\label{Sect::task}
\par
\renewcommand{\labelenumii}{\arabic{enumi}.\arabic{enumii}.}
Реализовать три варианта метода Гаусса с перестановками и научиться оценивать погрешность решения системы линейных уравнений для матриц произвольной размерности.

\section{Задачи}\label{Sect::task}
\par
\renewcommand{\labelenumii}{\arabic{enumi}.\arabic{enumii}.}
\begin{enumerate} 
    \item Реализовать метод Гаусса с перестановками по столбцам, по строкам, по столбцам и строкам одновременно для действительных квадратных матриц произвольной размерности n. 
    \item Для проверки работоспособности алгоритмов необходимо использовать алгоритм тестирования, который заключался в том, что мы заведомо определяем значения координат вектора x, данный вектор является решением уравнения A·х = b; вычисляем b путем прямого перемножения матрицы A на вектор x и далее производим поиск решения уравнения A·х = b тем или иным методом Гаусса,  получая x(числ), после чего производим сравнение полученного  x(числ) c заданным x, а также решением x(библ), полученным с использованием сторонней библиотеки выбранной студеном. При этом сравнение производится по Евклидовой норме разности вектора x-x(числ) и x-x(библ).
    \item На защите лабораторной работы студент должен показать умение оценивать погрешность вычислений в зависимости от выполнения  условия диагонального преобладания матрицы, умение сравнивать погрешности вычислений полученных методом Гаусса с перестановками по столбцам, по строкам, по столбцам и строкам одновременно. Понимать связь теории с практикой.
    \item Результат работы должен быть представлен в виде графиков зависимости абсолютной погрешности вычислений классическим методом Гаусса, методом Гаусса с перестановками по строкам, методом Гаусса с перестановками по столбцам, методом Гаусса с перестановками по столбцам и строкам, библиотечным методом от степени диагонального преобладания. Все графики должны быть построены на одной координатной плоскости. Напомним, что погрешность вычисления вектора x системы линейных алгебраических уравнений A·x=b тем или иным способом рассчитывается по Евклидовой норме разности точного решения и решения полученного соответствующим методом. Степень диагонального преобладания вычисляется, как максимальная разность по i между модулем диагонального элемента и суммы модулей вне диагональных элементов. Очевидно, что если значение степени диагонального преобладания положительна, то условие диагонального преобладания выполняется, в противном случае — не выполняется. Поэтому график должен быть построен как для отрицательных значений  степени диагонального преобладания, так и для положительных.
\end{enumerate}

\pagebreak
\section{Практическая реализация}
\par
В отчёте представлен вариант для квадратной матрицы размером 130х130 и амплитудой значений [-1000:1000]. В качестве заведомого известного вектора х был взят единичный вектор и посчитан ветор F. На основе матрицы А и вектора F с помощью различных методов был вычислен вектор Х. В результате посчитаны погрешности вычислений для классического метода Гаусса, метода строк, метода столбцов и комбинированного метода с различными коэффициентами диагонального преобладания.

Ниже представлен код программы:
\begin{lstlisting}
using Plots

n = 130
cup = 1000

function gauss_classic(A, F)
    for k in 1:n-1 
        for i in k+1:n
            dif = A[i,k]
            #A[i, :] ./= dif
            F[i] /= dif
            
            #A[i, :] .*= A[k, k]
            F[i] *= A[k, k]
            
            #A[i, :] .-= A[k, :]
            F[i] -= F[k]
            A[i, :] .= A[i,:] ./ dif .* A[k, k] .- A[k, :] 
            #F[i:n] .= F[i:n] ./ dif .*  A[k, k] .- F[k]
        end
   end

    X = ones(Float64, n)
    X[n] = F[n] / A[n,n]
    for i in n-1:-1:1
        dif = 0.0
        for j in i+1:n
            dif += A[i,j] * X[j]
        end
        X[i] = (F[i] - dif) / A[i, i]
    end
    return X
end

##строки
function gauss_strings(A, F)
    for k in 1:n-1
        max_i = k
        max_val = 0.0
        for i in k:n
            if (abs(max_val)) < abs(A[i,k]) 
                max_val = A[i,k]
                max_i = i
            end
        end
    
    
        A[k, :], A[max_i, :] = A[max_i, :], A[k, :]
        F[k], F[max_i] = F[max_i], F[k]

        
        for i in k+1:n
            dif = A[i,k]
            #A[i, :] ./= dif
            F[i] /= dif
            
            #A[i, :] .*= A[k, k]
            F[i] *= A[k, k]
            
            #A[i, :] .-= A[k, :]
            F[i] -= F[k]
            A[i, :] .= A[i,:] ./ dif .* A[k, k] .- A[k, :] 
            #F[i:n] .= F[i:n] ./ dif .*  A[k, k] .- F[k]
        end
   end

    X = ones(Float64, n)
    X[n] = F[n] / A[n,n]
    for i in n-1:-1:1
        dif = 0.0
        for j in i+1:n
            dif += A[i,j] * X[j]
        end
        X[i] = (F[i] - dif) / A[i, i]
    end
    return X
        
end

#столбцы
function gauss_cols(A, F)
    
    for k in 1:n-1
        max_i = k
        max_val = 0.0
        for i in k:n
            if (abs(max_val)) < abs(A[k,i]) 
                max_val = A[k,i]
                max_i = i
            end
        end

        A[:, k], A[:, max_i] = A[:, max_i], A[:, k]

        for i in k+1:n
            dif = A[i,k]
            #A[i, :] ./= dif
            F[i] /= dif
            
            #A[i, :] .*= A[k, k]
            F[i] *= A[k, k]
            
            #A[i, :] .-= A[k, :]
            F[i] -= F[k]
            A[i, :] .= A[i,:] ./ dif .* A[k, k] .- A[k, :] 
            #F[i:n] .= F[i:n] ./ dif .*  A[k, k] .- F[k]
        end        
    end

    X = ones(Float64, n)
    X[n] = F[n] / A[n,n]
    for i in n-1:-1:1
        dif = 0.0
        for j in i+1:n
            dif += A[i,j] * X[j]
        end
        X[i] = (F[i] - dif) / A[i, i]
    end
    return X
        
end


##стркои и столбцы
function gauss_sting_and_cols(A, F)
    for k in 1:n-1
        max_i = k
        max_val = 0.0
        for i in k:n
            if (abs(max_val)) < abs(A[i,k]) 
                max_val = A[i,k]
                max_i = i
            end
        end
    
        A[k, :], A[max_i, :] = A[max_i, :], A[k, :]
        F[k], F[max_i] = F[max_i], F[k]

        max_i = k
        max_val = 0.0
        for i in k:n
            if (abs(max_val)) < abs(A[k,i]) 
                max_val = A[k,i]
                max_i = i
            end
        end

        A[:, k], A[:, max_i] = A[:, max_i], A[:, k]

        
        for i in k+1:n
            dif = A[i,k]
            #A[i, :] ./= dif
            F[i] /= dif
            
            #A[i, :] .*= A[k, k]
            F[i] *= A[k, k]
            
            #A[i, :] .-= A[k, :]
            F[i] -= F[k]
            A[i, :] .= A[i,:] ./ dif .* A[k, k] .- A[k, :] 
            #F[i:n] .= F[i:n] ./ dif .*  A[k, k] .- F[k]
        end  
    end
    X = ones(Float64, n)
    X[n] = F[n] / A[n,n]
    for i in n-1:-1:1
        dif = 0.0
        for j in i+1:n
            dif += A[i,j] * X[j]
        end
        X[i] = (F[i] - dif) / A[i, i]
    end
    return X
        
end

function norm(x)
    sum = 0.0
    for element in x
        sum += element^2
    end
    return sqrt(sum)
end



X = ones(Float64, n)
#X .= 10
A = rand(-cup:cup, n, n)
A = Float64.(A)


F = [Float64(x) for x in (A * X)]

X0 = gauss_classic(copy(A), copy(F))
X1 = gauss_strings(copy(A), copy(F))
X2 = gauss_cols(copy(A), copy(F))
X3 = gauss_sting_and_cols(copy(A), copy(F))

println("Error in the library Gaussian method: ", norm(X - X00))
println("Error in the classical Gaussian method: ", norm(X - X0))
println("Error in the string method: ", norm(X - X1))
println("Error in the column method: ", norm(X - X2))
println("Error in the combined method: ", norm(X - X3))

gauss_classic_array = Vector{Float64}()
gauss_strings_array = Vector{Float64}()
gauss_cols_array = Vector{Float64}()
gauss_sting_and_cols_array = Vector{Float64}()
diag_coefs = Vector{Float64}()


diagonal_increase_coef = -15.0
while diagonal_increase_coef < 15.0
    push!(diag_coefs, diagonal_increase_coef)
    cool_A = copy(A[1:n, 1:n])

    for i in 1:n
        if cool_A[i, i] == 0.0
            cool_A[i, i] += 1
        end
        cool_A[i, i] += 2^diagonal_increase_coef * sign(cool_A[i, i] )
    end

    F1 = [Float64(x) for x in (cool_A * X)]

    push!(gauss_classic_array, norm(X - gauss_classic(copy(cool_A), copy(F1))))
    push!(gauss_strings_array, norm(X - gauss_strings(copy(cool_A), copy(F1))))
    push!(gauss_cols_array, norm(X - gauss_cols(copy(cool_A), copy(F1))))
    push!(gauss_sting_and_cols_array, norm(X - gauss_sting_and_cols(copy(cool_A), copy(F1))))
    diagonal_increase_coef += 1

end



plot(diag_coefs, gauss_classic_array, label="Classic", xlabel="Diagonal coef", ylabel="Error")
plot!(diag_coefs, gauss_strings_array, label="Strings", color=:red)
plot!(diag_coefs, gauss_cols_array, label="Columns", color=:yellow)
plot!(diag_coefs, gauss_sting_and_cols_array, label="Combination", color=:green)
\end{lstlisting}

\section{Результаты}
\par
Результат для матрицы с нулевым коэффициентом диагонального преобладания:
\begin{center}
\includegraphics[scale=0.75]{g1}
\par
\end{center}

Результаты вычислений различных методов с наращиванием коэффициента диагонального преобладания:
\begin{center}
\includegraphics[scale=0.75]{g2}
\par
\end{center}

\section{Выводы}
В ходе выполнения лаборатороной работы были реализованы три варианта метода Гаусса с перестановками и посчитаны погрешности решения системы линейных уравнений для матриц произвольной размерности. В сравнении с остальными методами классический метод Гаусса оказался гараздо менее точным. Наиболее надёжным же оказался комбинарованный метод.
\end{document}
